\chapter{Experimental Results}
\label{chapter_experimentalResults}

\section{Temporal Service-Value Decay Experiment}
\label{sec:temporal_decay_experiment}

Before running the full responder-UAV planning experiment, an initial
calibration experiment was performed to examine the time-dependent service
value function used in the PGBM objective. The purpose of this experiment was
to identify a simple severity-aware priority function that decreases with
service delay while preserving the ordering between tasks with different
operational severity scores.

This experiment was necessary because the service-value function directly
controls how the optimization model prioritizes tasks. If only the fixed
severity score is used, the model cannot represent the loss of benefit caused
by delayed assistance. On the other hand, if the decay function is too
aggressive for high-severity tasks, the model may unintentionally allow a
lower-severity task to become more valuable than a higher-severity task after
some delay. The selected function therefore needs to represent time pressure
without violating the intended severity ordering.

The experiment considered seven operational severity levels from 0.3 to 0.9
and evaluated service value over a 120-minute post-detection period. A
reference policy table was first prepared to represent a plausible operational
assumption: tasks with higher severity start with higher value, and every task
loses value as service is delayed. The reference values were not treated as
measured medical survival probabilities; they were used only as a calibration
guide for selecting a service-value function.

\begin{table}[htbp]
\centering
\caption{Reference policy values used for service-value calibration.}
\label{tab:reference_policy_values}
\begin{tabular}{ccccc}
\hline
\textbf{Severity} & \textbf{0 min} & \textbf{30 min} & \textbf{60 min} & \textbf{120 min}\\
\hline
0.9 & 0.90 & 0.65 & 0.40 & 0.15\\
0.8 & 0.80 & 0.60 & 0.38 & 0.15\\
0.7 & 0.70 & 0.53 & 0.35 & 0.14\\
0.6 & 0.60 & 0.46 & 0.31 & 0.13\\
0.5 & 0.50 & 0.38 & 0.27 & 0.12\\
0.4 & 0.40 & 0.31 & 0.23 & 0.11\\
0.3 & 0.30 & 0.25 & 0.19 & 0.10\\
\hline
\end{tabular}
\end{table}

Table~\ref{tab:reference_policy_values} shows the reduced version of the
reference policy used in the experiment. The complete notebook also included
intermediate values at 10 and 90 minutes. The table reflects the intended
behavior of the service-value model: value decreases with delay, while
higher-severity tasks remain higher in value than lower-severity tasks at the
same delay.

\begin{figure}[htbp]
    \centering
    \includegraphics[width=0.82\textwidth]{reference_policy_service_value.png}
    \caption{Reference policy curves used to guide service-value calibration.}
    \label{fig:reference_policy_service_value}
\end{figure}

Figure~\ref{fig:reference_policy_service_value} visualizes the same reference
policy values. The graph makes the intended trend clearer than the table alone:
each severity level starts from its assigned operational severity score, and
the reference value decreases as service delay increases.

Four candidate functions were compared: a common exponential decay function,
a severity-dependent exponential decay function, a severity-dependent
hyperbolic decay function, and a severity-dependent linear decay function.
For each candidate, the decay parameter was fitted by minimizing the mean
squared error between the candidate curve and the reference policy values.
After fitting, each function was checked to ensure that service value never
increases with delay and that a higher-severity task does not fall below a
lower-severity task at the same delay.

\begin{table}[htbp]
\centering
\caption{Candidate service-value functions evaluated in the experiment.}
\label{tab:candidate_decay_functions}
\begin{tabular}{p{0.32\textwidth}p{0.46\textwidth}}
\hline
\textbf{Candidate} & \textbf{Function Form}\\
\hline
Common exponential & $V(s,t)=s e^{-\lambda t}$\\
Severity exponential & $V(s,t)=s e^{-\lambda s t}$\\
Severity hyperbolic & $V(s,t)=\dfrac{s}{1+\lambda s t}$\\
Severity linear & $V(s,t)=\max(0,s(1-\lambda s t))$\\
\hline
\end{tabular}
\end{table}

Table~\ref{tab:candidate_decay_functions} lists the four candidate forms.
Here, $s$ denotes the operational severity score and $t$ denotes the delay
after task detection. The common exponential function applies the same decay
rate to every task and uses the severity score as the initial value. The
severity-dependent functions allow the rate of value loss to change with the
severity level.

\begin{figure}[htbp]
    \centering
    \includegraphics[width=0.92\textwidth]{candidate_service_value_curves.png}
    \caption{Fitted service-value curves for the four candidate functions.}
    \label{fig:candidate_service_value_curves}
\end{figure}

Figure~\ref{fig:candidate_service_value_curves} shows how the four fitted
functions behave across the same delay window. The common exponential keeps
the curves separated, while the severity-dependent exponential and linear
forms allow curves to move closer or cross over time.

\begin{table}[htbp]
\centering
\caption{Comparison of candidate service-value decay functions.}
\label{tab:service_value_decay_comparison}
\begin{tabular}{p{0.34\textwidth}ccc}
\hline
\textbf{Candidate Function} & \textbf{Fitted $\lambda$} & \textbf{MSE} & \textbf{Ordering}\\
\hline
Common exponential & 0.011997 & 0.000892 & Preserved\\
Severity exponential & 0.017029 & 0.001259 & Not preserved\\
Severity hyperbolic & 0.026040 & 0.002742 & Preserved\\
Severity linear & 0.010213 & 0.003129 & Not preserved\\
\hline
\end{tabular}
\end{table}

Table~\ref{tab:service_value_decay_comparison} shows that the common
exponential function produced the smallest fitting error while also satisfying
both safety checks. The severity-dependent exponential and linear functions
were rejected because they could reverse the intended priority ordering: a
lower-severity task could retain more value than a higher-severity task after
some delay. The severity-dependent hyperbolic function preserved the ordering,
but its fitting error was larger than the common exponential alternative.

\begin{figure}[htbp]
    \centering
    \includegraphics[width=0.82\textwidth]{candidate_function_mse_comparison.png}
    \caption{Fitting error comparison for the candidate service-value functions.}
    \label{fig:candidate_function_mse_comparison}
\end{figure}

Figure~\ref{fig:candidate_function_mse_comparison} summarizes the model
selection result. The common exponential candidate has the lowest fitting
error and preserves the severity ordering, which makes it the most suitable
choice among the tested functions.

The result is important for the proposed PGBM formulation because the selected
function must be useful inside an optimization objective, not only visually
reasonable in a plot. The common exponential function is smooth, monotone,
easy to interpret, and controlled by a single decay parameter. The fitted
decay rate also gives a direct way to tune how strongly delay reduces the
operational value of a task. For these reasons, it was selected as the most
appropriate service-value function for the current model.

\begin{figure}[htbp]
    \centering
    \includegraphics[width=0.86\textwidth]{service_value_decay_experiment.png}
    \caption{Fitted common exponential service-value decay curves. Markers
    indicate the reference policy values used for calibration.}
    \label{fig:service_value_decay_experiment}
\end{figure}

Figure~\ref{fig:service_value_decay_experiment} illustrates the selected
common exponential model. The fitted curves decrease smoothly with delay, and
the separation between severity levels is preserved throughout the 120-minute
window. For example, with the fitted decay rate, a task with severity 0.9 has
service value approximately 0.628 after 30 minutes and 0.438 after 60 minutes,
whereas a task with severity 0.6 has value approximately 0.419 after 30
minutes and 0.292 after 60 minutes. Thus, delay reduces the value of all
tasks, but the relative priority implied by the operational severity score
remains stable.

\section{Results Analysis}
\label{result_analysis}

The experiment supports the use of the common exponential service-value model
already adopted in the proposed formulation:
\[
v_i(C_i)=s_i e^{-\lambda(C_i-t_i)}.
\]
This function is simple, monotone, and directly interpretable. The severity
score $s_i$ determines the initial operational priority of the task, while the
decay parameter $\lambda$ controls how quickly the value decreases with
completion delay. Since the same decay rate is applied to all tasks, the model
avoids unintended priority reversals and expresses delay sensitivity through a
single tunable parameter.

The selected function should still be interpreted as an operational priority
model, not as a validated medical survival model. For the current stage of the
thesis, it provides a suitable balance between mathematical simplicity,
interpretability, and consistency with the golden-hour motivation.

\section{Summary of the Experimental Results}
\label{experimentsummary}

The temporal service-value experiment indicates that the common exponential
decay function is the preferred service-value model for PGBM. It achieved the
lowest mean squared error among the tested candidates and preserved the
priority ordering between operational severity levels. Therefore, this thesis
uses the common exponential service-value function in the optimization
objective. The result justifies the current formulation choice before the full
task-allocation, 3D routing, payload, energy, and recourse experiments are
implemented.
